% =====================================================================
% Reviewer gbSA -- Rating 3 (Borderline reject), Confidence 3
% Owns Rebuttal Tables 5 and 6.
% Cross-references Rebuttal Tables 1, 2, 3, 4, 7, 8, 9. Reprints none of them.
% =====================================================================

\section*{Reviewer gbSA}

\noindent\textit{Rating 3 (Borderline reject), Confidence 3. Four weaknesses, four questions,
plus a limitations entry of the reviewer's own.}

\vspace{0.4em}
\noindent
This is the most demanding of the three reviews, and the one that changed the most in our work:
a rebuttal-period experiment supports Weakness 2 against our own earlier reading, and that
reversal is reported in full below.

% ---------------------------------------------------------------------
\subsection*{Weaknesses}
% ---------------------------------------------------------------------

\begin{itemize}
\item[\textbf{W1.}] \itshape ``The biggest conceptual issue is that the title and framing still
feel overly anthropomorphic. Even though the paper says that `addiction-like' is only a
behavioral label, the title and the frequent use of clinical terminology may still lead readers
to think the paper is asking whether LLMs can literally become `addicted.' ''
\upshape
\plan{임상 용어를 제목에서 제거하고, 초록 첫 문장을 조작적 정의로 교체하며, 본문 전체 임상 어휘를 감사한 새 제목을 제시한다.}

\item[\textbf{W2.}] \itshape ``In this paper, the LLMs are placed into several role-playing
settings and asked to make choices. In this context, continuing to play, setting goals, or
trying to make profits may simply reflect instruction following, role-play priors, or
misunderstanding the intended task objective. It is not clear that these behaviors reflect
stable risk preferences or a general mechanism induced by autonomy.''
\upshape
\plan{페르소나 서문이 Gemini에서 파산 엔드포인트를 실제로 움직인다는 신규 요인설계 결과와 거절-내용 감사 결과를 자기반증으로 보고하고, 안정적 위험선호(trait) 해석을 철회하며, 순수 지시추종 설명이 수용하지 못하는 목표 상향 논거 하나만 남긴다.}

\item[\textbf{W3.}] \itshape ``Most of the main figures do not include confidence intervals or
error bars. Many results are reported only with `p < .05,' without giving enough detail about
effect-size uncertainty.''
\upshape
\plan{모든 셀에 n과 95\% Wilson 구간을 붙이고, 아직 구간이 없는 부분(arm 간 $\Delta$)과 사전등록 판정규칙 2건의 실패를 포함한 통계적 결함 5종을 그대로 공개한다.}

\item[\textbf{W4.}] \itshape ``The comparisons in the fixed-vs-variable slot machine experiment
and the investment choice experiment are not fully fair. For example, variable betting does not
only introduce `freedom of choice'. It also changes the action space, the available strategies,
the game length, and stopping behavior.''
\upshape
\plan{일치 상한이 라운드당 최대 스테이크만 동일화해 `범위 확장' 설명을 반증할 뿐 누적 노출·행동공간·전략공간은 동일화하지 못함을 인정하고, 게임 길이 조건부 분석이 사전등록상 불가능한 이유를 설명하며, 투자선택 과제의 공정성 기반 추론은 철회한다.}
\end{itemize}

% ---------------------------------------------------------------------
\subsection*{Questions}
% ---------------------------------------------------------------------

\begin{itemize}
\item[\textbf{Q1.}] \itshape ``Why is the matched-cap ablation only run on one GPT-4o-family
model? I would suggest testing this on several other models as well.''
\upshape
\plan{6 모델 $\times$ 4 상한 $\times$ 2 모드, 셀당 200게임, 48/48 셀 완주로 확장했고 결과는 Reviewer KuK5 절의 Rebuttal Table 1을 참조하도록 상호참조한다.}

\item[\textbf{Q2.}] \itshape ``For the keyword-based cognitive distortion detection, what are the
keyword lists, annotation rules, and human validation results?''
\upshape
\plan{키워드 목록(부록 A, 동결 파일 해시 포함)과 주석 규칙(4단계 척도)은 전부 제시하고, 인간 검증 결과는 라벨을 한 건도 수집하지 않았으므로 존재하지 않는다고 명시한다.}

\item[\textbf{Q3.}] \itshape ``In the slot machine task, does the model explicitly know that
stopping immediately is the EV-optimal action? If not, continuing to play may just mean that the
model interprets the task as `you are supposed to play this game.' Did the authors test an
explicit rationality instruction or a more decision-theoretic framing?''
\upshape
\plan{명시적 합리성 지시를 페르소나·자율성과 교차한 요인설계를 실행해 Rebuttal Table 5·6으로 보고하고, 파산이 대부분 셀에서 0이지만 전부는 아니라는 점과 GPT-4.1-mini의 상시 시스템 메시지를 포함한 한계 4종을 붙인다.}

\item[\textbf{Q4.}] \itshape ``In the internal representation analysis, why use SAE top-200
features plus Ridge regression, instead of directly using logit or choice-probability controls,
or a simpler behavioral-state baseline?''
\upshape
\plan{SAE 상위-k + Ridge를 택한 이유가 통계적 우월성이 아니라 가독성이었음을 그대로 밝히고, 실제로 실행한 통제(SAE 없는 raw-ridge)와 실행하지 않은 통제(중첩 행동상태 베이스라인)를 구분해 후자는 미해결 약점으로 남긴다.}
\end{itemize}

% ---------------------------------------------------------------------
\subsection*{Limitations (the reviewer's own entry)}
% ---------------------------------------------------------------------

\begin{itemize}
\item[] \itshape ``The slot machine, investment choice, and mystery wheel tasks are all highly
artificial negative-EV games. They cannot directly represent the risk behavior of real financial
agents, planning agents, or tool-using agents. Overall, I think the results are better understood
as exploratory monitoring signals rather than as evidence of an underlying mechanism.''
\upshape
\plan{이 진술을 논문 자신의 입장으로 채택하되 범위는 내부상태 층으로 한정하고, 기여를 행동층·내부층 2단으로 재진술해 증거 지위가 부록을 읽지 않아도 보이게 한다.}
\end{itemize}

\redrule

% =====================================================================
\subsection*{\color{red!60!black}Author Response}
% =====================================================================

We thank you for writing each concern as a testable statement: one of them --- Weakness 2 ---
is now supported by an experiment we ran during this rebuttal, against our own earlier reading.
Your points press on one joint: whether the effects we report belong to the models or to our
apparatus --- the framing, the prompts, the comparison structure and the statistics. This
section is written to be read once, in order, without the paper open; here is every answer in
brief.

\begin{itemize}
\item \textbf{Is the title and framing overly anthropomorphic? (W1)} Yes. The clinical term
leaves the title, the abstract opens with the operational definition, and clinical vocabulary is
audited throughout.
\item \textbf{Could the behaviour be role-play priors or instruction following? (W2)} Partly,
and our own new data support you against us: the persona preamble alone moves Gemini's
pre-registered bankruptcy endpoint (fixed 0 to 12, variable 4 to 32, per 100 games). We
withdraw the stable-risk-preference reading; one observation still resists a pure
instruction-following account.
\item \textbf{Where are the confidence intervals? (W3)} Added. Every matched-cap cell now
carries a 95\% Wilson interval on a stated $n$; the arm-to-arm $\Delta$ still lacks one, due by
3~August, and both pre-registered decision rules failed --- disclosed rather than repaired.
\item \textbf{Is fixed-versus-variable a fair comparison? (W4)} Not fully, exactly as you say.
The matched cap removes only the range-expansion reading, and in one model strongly undercuts it: LLaMA's
smaller-stake arm is the one that ruins (81.5\% vs 3.0\%) --- while cumulative exposure and the
action space stay unequalised, and we withdraw the investment-task fairness inference.
\item \textbf{Why was the matched cap run on only one model? (Q1)} It no longer is: 6 models
$\times$ 4 caps $\times$ 2 modes at $n = 200$ per cell, all 48 cells complete (Rebuttal Table 1,
Reviewer KuK5's section). Decisive in LLaMA, weak in Gemini, untestable in the four models at
0.0\% in both arms.
\item \textbf{Keyword lists, annotation rules, human validation? (Q2)} Lists: Appendix A, with
the frozen file's hash. Rules: the four-level scale printed below. Validation: none yet --- the
blinded instrument is deployed and labels follow by 3~August.
\item \textbf{Does the model know stopping is EV-optimal? (Q3)} Tested. An explicit rationality
instruction removes almost all play in the three models that play most ($-91$ to $-100$ pp
participation), though it is not a complete off switch and cannot yet separate information from
endorsement.
\item \textbf{Why SAE top-200 + Ridge rather than simpler controls? (Q4)} Legibility, not
statistical superiority. The behavioural-state baseline you name has not been run; it stays an
open weakness, registered for 3~August.
\item \textbf{Your limitations entry.} Adopted, with one scope line: the internal-state results
are claimed as exploratory monitoring signals, while the behavioural layer keeps its status as a
pre-registered-endpoint experimental effect.
\end{itemize}

If any part of our response falls short, we would be glad to take it further during the
discussion period.

% ---------------------------------------------------------------------
\subsubsection*{\textbf{[Weakness 1]} Anthropomorphic title and framing}
% ---------------------------------------------------------------------

\textbf{Verdict: accepted.} The title announces a claim the paper never makes, and a disclaimer
that arrives after the first impression cannot repair it. The proposed title is:

\begin{center}
\begin{minipage}{0.86\linewidth}
\itshape Autonomy and Gambling-Like Risk-Taking in Large Language Models: Behavioural Evidence
and Conditional Internal-State Readouts.
\end{minipage}
\end{center}

\noindent
Three edits follow: the clinical term leaves the title; the abstract now opens with the
operational definition instead of placing it in the second sentence; and clinical vocabulary is
replaced by behavioural description wherever the clinical term was unnecessary. Codebook
construct names keep their published names, which tie them to their source instruments, and are
identified as instrument names, not diagnoses, on first use (lineage: Rebuttal Table 8,
Reviewer a3Zu's section). The new title says ``readouts'', not ``evidence'', because
``evidence'' would suggest a mechanism our causal battery does not establish, and
``conditional'' records that the readout strengthens under autonomy; Reviewer KuK5 raises the
same issue.

However, we think the clinical measurement framework itself should stay. The paper does not ask
whether a model can literally become addicted; it asks when an LLM's sequential decisions
become irrational, and clinical gambling research is what makes that observable round by round.
We intend to reflect this in the camera-ready.

% ---------------------------------------------------------------------
\subsubsection*{\textbf{[Weakness 2]} Role-play priors, instruction following, task misunderstanding}
% ---------------------------------------------------------------------

\textbf{Verdict: accepted, and our own new data support you against us.} We crossed the persona
preamble used verbatim by the paper's open-weight runs, an explicit rationality instruction,
and autonomy over stake size, at cap \$70 with $n = 100$ games per cell (32 of 44 cells
complete; Rebuttal Table 5 under Question 3).

\vspace{0.3em}
\noindent\textbf{We correct an earlier claim of ours.} We introduced the preamble to prevent
safety refusals, treated it as behaviourally neutral, and earlier said it did essentially
nothing. Gemini-2.5-Flash contradicts that: the preamble raises fixed-arm participation from 8
to 53 of 100 games, with 12 bankruptcies where there were none, and raises variable-arm mean
rounds from 2.25 to 8.43 with bankruptcies from 4 to 32. Role-play framing alone moves the
pre-registered primary endpoint in at least one model (Appendix B10 records the withdrawn
claim).

\vspace{0.3em}
\noindent\textbf{The safety-filter explanation does not survive its own audit.} The persona does
remove safety-style declining entirely in all four API models. But safety language appears in
only 0--16\% of no-persona refusals, against expected-value reasoning in 78--100\%, so removing
Gemini's 12\% sliver cannot account for its 45-point participation rise. Stronger still:
Gemini's variable arm is at ceiling in both cells, yet rounds and bankruptcies rise over the
same 100 games --- a device that only removes start-up refusals cannot alter play after every
game has begun. What accounts for the remainder is not established; the most defensible reading,
labelled as a reading, is that simulation framing makes losses feel costless.

\vspace{0.3em}
\noindent\textbf{The effect is model-dependent, and autonomy remains the larger lever.} The
same preamble moves variable-arm participation by only $+1.0$ to $+2.0$ pp (percentage points)
in the two GPT models, against an autonomy effect of $+92.0$ to $+97.0$ pp in the same cells
(Rebuttal Table 6). Two limits: the casino framing is never removed, and the preamble bundles
three elements --- a decomposition arm removing its compliance sentence is registered and unrun
(Rebuttal Table 9), reported by 3~August.

\vspace{0.3em}
\noindent\textbf{However, two observations resist a pure instruction-following account.}
Language: after deleting every expression the goal instruction itself could have supplied, the
goal minus no-goal contrast stays positive in 6 of 6 models at $+3.1$ to $+41.9$ pp, carried by
frames nothing in the instruction asks a model to express (Rebuttal Table 7, Reviewer a3Zu's
section; a language measure, not a behavioural one). Behaviour: raising a goal \emph{after}
meeting it is not something any instruction requests, and under that stricter definition the
goal arm stays at $2.24\times$ the no-goal arm on the sample where the instrument is symmetric,
positive in all four models where the contrast is defined (full analysis and its two instrument
defects: Reviewer a3Zu's Question 1).

\vspace{0.3em}
\noindent\textbf{We withdraw the trait reading of ``stable risk preferences'', because you are
right.} LLaMA's fixed-arm bankruptcy is not monotone in the cap (0.0 / 0.5 / 13.0 / 3.0\% at
caps \$10--\$70; Rebuttal Table 1), driven by willingness to re-bet after a first loss falling
from 57\% to 6\% across caps. We now report these indices as condition-dependent policies, and
Question 3 below addresses task misunderstanding directly. We intend to reflect this in the
camera-ready.

\plan{제목·초록·본문에서 trait 어휘를 제거하고, 지표를 조건의존 정책으로 재정의하며, 페르소나 서문을 관심 조작 변수로 승격해 요인설계·거절내용 감사 결과와 함께 본문에 싣는다.}

% ---------------------------------------------------------------------
\subsubsection*{\textbf{[Weakness 3]} Missing confidence intervals, and ``p < .05'' without effect-size uncertainty}
% ---------------------------------------------------------------------

\textbf{Verdict: accepted, and one half of it is still open.} Recomputing every cell with an
interval is also how we found that both pre-registered decision rules had failed.

\vspace{0.3em}
\noindent\textbf{What we added.} Every matched-cap cell now carries a 95\% Wilson interval and a
stated $n$ (48 of 48 cells at $n = 200$; Rebuttal Table 1), the revision adds intervals to
every main figure, and factorial cells are labelled as raw counts.
\textbf{What is still missing.} The $\Delta$ between arms has no interval and no standardised
effect size; a recomputation clustering by model on the probability scale is registered but
unrun (Rebuttal Table 9), reported by 3~August.

\vspace{0.3em}
\noindent\textbf{Five further disclosures follow.}

\textit{(i) The primary pre-registered rule is ill-posed for this design.} It asks whether the
lower 2.5\% posterior quantile of the primary coefficient exceeds zero, but the fixed arm sits
at 0.0\% in every tabulated cell of four of the six models and never exceeds 13.0\%, so no data
this design produces could reject the rule. Passing it is not evidence.

\textit{(ii) Our analysis output wrongly recorded that rule as passing}, on a pooled resampling
interval ($+14.25$ pp [13.25, 15.25]) the frozen configuration neither specifies nor names as a
fallback, while the registered mixed-effects model produced non-finite cluster-robust errors.
The record is wrong; this response corrects it.

\textit{(iii) The qualitative secondary rule is met by 0 of 6 models.} The exposure clause
passes in 5 of 6; the stake-size clause fails in 5 of 6, because models given discretion wager
22--46\% of the cap, not more than half. That failure is informative: we expected free models
to bet bigger, and they instead bet small and kept playing.

\textit{(iv) Code and configuration state the secondary rule differently.} The code implements
a third rule, which also fails (2 of 6), while labelling its output with the configuration's
clause text. All three readings fail; we correct the audit trail instead of choosing among
them.

\textit{(v) Four provenance and execution defects} (Appendix B1, B5, B7 and B8, each with
measured impact and prevention): the pre-registration freeze date postdates a registered model's
removal from the vendor API; the configuration names \texttt{claude-3-5-haiku-20241022} while
every run used \texttt{claude-haiku-4-5-20251001}; the two harnesses use different fixed-arm
round caps (50 and 100); and a stake-execution defect made fixed cells at caps \$30--\$70
execute \$10 instead of the cap, so those 18 cells were re-run and every such cell in Rebuttal
Table 1 is post-fix. The prompt text itself is clean: byte-identical first-round prompts across
all six models (SHA-256 prefix \texttt{704a35b8e22f34de}, 335 characters).

\vspace{0.3em}
\noindent\textbf{One measurement defect changes a number in this letter.} Our decision parser
took the first ``final decision'' match and tested bet-tokens before stop-tokens, so a response
arguing for walking away and ending with a stop could be scored as a wager. Re-parsing every
stored response under a corrected rule flips at most 0.3\% of adjudicable decisions in any
corpus (Rebuttal Table 4, appendix, with denominators) --- but all 14 flips in the matched-cap
re-run are Claude-Haiku, so that model's small fixed-arm participation counts are essentially
artefacts, annotated wherever they appear. We intend to reflect this in the camera-ready.

\plan{모든 본문 그림에 n과 95\% Wilson 구간을 넣고, 판정규칙 3종의 실패와 실행·출처 결함을 부록의 자기공개 항목으로 고정하며, arm 간 $\Delta$ 구간 재계산은 2차 응답으로 등록한다.}

% ---------------------------------------------------------------------
\subsubsection*{\textbf{[Weakness 4]} Fixed versus variable is not a fair comparison}
% ---------------------------------------------------------------------

\textbf{Verdict: accepted in part, and this is the objection we take most seriously.} Variable
betting changes far more than freedom of choice --- action space, strategies, game length and
stopping all move with it --- so we state which part the matched cap removes and which it
leaves untouched.

\vspace{0.3em}
\noindent\textbf{What the matched cap fixes.} It equalises the maximum stake per round,
removing the claim that the variable condition simply permits larger bets --- and the data
undercut that claim where the endpoint is informative. At cap \$70, LLaMA's fixed arm
executes a mean wager of \$68.4 against the variable arm's \$32.1, yet the larger-stake arm
produces 3.0\% bankruptcy against 81.5\%: range expansion predicts the opposite ordering
(Rebuttal Table 2, Reviewer KuK5's section). Gemini shows a weak contrast in the same
direction; the other four models are at 0.0\% in both arms, so their ordering cannot invert.

\vspace{0.3em}
\noindent\textbf{What it does not fix, exactly as you say.} Cumulative exposure: models play
longer when they choose their stakes (Rebuttal Table 3; its Claude-Haiku fixed-arm column is a
parsing artefact, audited in Rebuttal Table 4), and we removed the sentences implying we had
separated discretion from the exposure it produces. Action space: the cap equalises the range
ceiling, not the number of available actions, and only the variable arm can implement policies
that require additional stakes, escalation after a loss being the clearest. On naming: the
submitted abstract already describes the manipulation operationally --- ``letting the model
choose its own bet size'' --- while the word ``freedom'' appears in \S1, \S3 and \S5, and the
camera-ready replaces it there with \emph{discretion over stake size}, counting the effect on
the action set as part of the manipulation.

\vspace{0.3em}
\noindent\textbf{Why we do not condition on game length.} Our frozen configuration contains no
length-conditioned analysis, so adding one now would be exactly the unregistered post hoc move
that item (ii) of Weakness 3 discloses against us; game length and stopping are outcomes of the
manipulation, so conditioning on them conditions on a post-treatment variable; and length and
ruin share an upstream cause, so splitting games by length would manufacture association. The
appropriate repair is ex-ante --- the same round budget in both arms, per-round risk of ruin at
matched stakes --- registered and unrun (Rebuttal Table 9), reported by 3~August.

\vspace{0.3em}
\noindent\textbf{We did not address the investment-choice task and should have.} We have no
matched-cap analogue for it and built none during the rebuttal, so we withdraw the
fairness-based inference: that task's fixed-versus-free contrast is now descriptive, and the
claim about discretion over stake size rests only on the slot-machine grid.

\vspace{0.3em}
\noindent\textbf{We now treat refusal to play as the rational baseline, which is your framing.}
At cap \$70 most models refuse the fixed-arm game, wagering in 0.0--7.5\% of games
(Claude-Haiku's figure is itself mostly misparsed stops; Rebuttal Table 3), while LLaMA wagers
in 66.0\% of fixed games yet only 3.0\% end in ruin --- against 81.5\% when it sets its own
stake. Discretion over stake size converts correct refusal into sustained play, and in one
model into ruin in 81.5\% of games; we rewrote the paper around that claim, and we intend to
reflect this in the camera-ready.

\plan{본문(\S1·\S3·\S5)의 조작 명칭을 `freedom'에서 `discretion over stake size'로 바꾸고(초록은 이미 조작적 표현이라 수정 불요), 누적 노출·행동공간 미통제를 한계로 명시하며, 동일 라운드 예산 설계를 2차 응답용으로 등록하고, 투자선택 과제의 공정성 기반 추론은 삭제한다.}

% ---------------------------------------------------------------------
\subsubsection*{\textbf{[Question 1]} Why is the matched-cap ablation run on only one model?}
% ---------------------------------------------------------------------

\textbf{It is no longer limited to one model; Reviewer KuK5's Question 1 contains our full
answer, and we keep only the conclusion here.} We ran 6 models $\times$ 4 caps $\times$ 2 modes
at $n = 200$ per cell, completed all 48 cells, and withdraw the claim that the matched-cap
dissociation generalises across the panel: it is decisive in LLaMA, weak in Gemini, and
untestable in the other four, whose bankruptcy is 0.0\% in both arms --- a silence that fails
to test the contrast rather than refuting it. The floors belong to the base prompt condition,
not to the models: under the paper's richer prompt conditions the same four models move far
above the floor (the two designs and those figures: KuK5's Question 1 and Appendix B9). We
intend to reflect this in the camera-ready.

\plan{6모델 매칭 상한 격자를 본문에 넣고, 6모델 일반화 주장을 철회하며, 바닥 효과가 API 드리프트가 아니라 base 조건의 성질이라는 정정을 부록 자기공개 항목으로 고정한다.}

% ---------------------------------------------------------------------
\subsubsection*{\textbf{[Question 2]} Keyword lists, annotation rules, and human validation results}
% ---------------------------------------------------------------------

\textbf{We can give you the first two, but not the third.} Appendix A contains the keyword
lists. The annotation rules are stated below. Human validation results do not exist yet, as no
labels are collected.

\vspace{0.3em}
\noindent\textbf{Keyword lists.} Appendix A prints every expression in all four constructs, with
its mapping to the submitted manuscript's frames. The frozen artefact is
\texttt{convergent\_codebook.FROZEN.py}, whose SHA-256 is

\begin{center}
\footnotesize\texttt{7d16e30d7d69284ae37493cf61fffcfb9db0b80ef69d3689fb1d13dfbb5e69d7}
\end{center}

\noindent
The four constructs recur across the validated instruments reviewed in Goodie and Fortune
(2013); we wrote the expressions ourselves after searching for, and not finding, a validated
public lexicon for gambling-specific distortions in free text (lineage: Rebuttal Table 8,
Reviewer a3Zu's section).

\vspace{0.3em}
\noindent\textbf{Annotation rules.} One trace, one construct, one question --- does this
response use [construct] as grounds for its next action? --- on a four-level scale.

\begin{enumerate}[leftmargin=1.8em]
\item \textbf{Uses as grounds.} The response invokes the construct and acts on it.
\item \textbf{Mentions but rejects.} The response raises the construct and declines to act on it.
\item \textbf{Unrelated.} The construct does not appear in any load-bearing role.
\item \textbf{Cannot tell.} The trace is too short, too truncated or too ambiguous to adjudicate.
\end{enumerate}

\noindent
A model that discusses a distortion in order to reject it receives level 2, not level 1; a
keyword match cannot make that distinction, which is why human review exists. Our regexes
already show two false positives: ``stopping now is the smart decision'', a rational refusal,
scores as self-serving bias, and illusion-of-control patterns misfire in the variable arm,
where the model genuinely controls its stake.

\vspace{0.3em}
\noindent\textbf{Human validation results do not exist yet.} The instrument is deployed and the
coders are blinded (design and pre-fixed decision rule: Reviewer a3Zu's Weakness 1), but no
labels are collected and the non-author coder is not yet recruited. We will complete it and
report by 3~August, regardless of whether the outcome favours us (Rebuttal Table 9).

\vspace{0.3em}
\noindent\textbf{However, two claims remain testable without human labels.} Across 19,200 games,
the goal versus no-goal contrast is positive in 6 of 6 models under the paper's instrument with
its window scoping restored, and stays positive in 6 of 6 after deleting every goal-coupled
expression, at $+3.1$ to $+41.9$ pp (LLaMA is a boundary case). Rebuttal Table 7 in Reviewer
a3Zu's section gives the figures and the three limits that attach. We intend to reflect this in
the camera-ready.

\plan{키워드 표현·동결 해시·4단계 주석 규칙을 부록 A로 공개하고, 키워드 척도를 도구(instrument)에서 특이도 미확립 베이스라인으로 강등하며, 인간 라벨은 2차 응답에서 결과와 무관하게 보고한다.}

% ---------------------------------------------------------------------
\subsubsection*{\textbf{[Question 3]} Does the model know stopping is EV-optimal, and did you test a rationality instruction?}
% ---------------------------------------------------------------------

\textbf{Verdict: accepted, and tested.} We crossed an explicit rationality instruction ---
stating the per-round expected loss, that stopping immediately maximises expected value, and
that the model may stop at any time --- with the persona preamble and autonomy at cap \$70,
$n = 100$ games per cell, 32 of 44 cells complete. Rebuttal Table 5 gives one row per completed
cell.

\vspace{0.6em}
\noindent\textbf{Rebuttal Table 5.} Framing $\times$ rationality factorial at cap \$70, n = 100
games per cell, all 32 completed cells of 44. Cells are raw counts, not rates with intervals.
ROLE is the persona preamble; RAT is the rationality instruction.

\vspace{0.3em}
{\small
\begin{tabular}{lllccrc}
\toprule
\textbf{Model} & \textbf{Mode} & \textbf{ROLE} & \textbf{RAT} & \textbf{$\ge$1 wager} & \textbf{Mean rounds} & \textbf{Bankrupt} \\
\midrule
Claude-Haiku & fixed    & none & 0 & 5/100  & 0.06 & 0 \\
             & fixed    & none & 1 & 2/100  & 0.02 & 0 \\
             & fixed    & role & 0 & 4/100  & 0.04 & 0 \\
             & fixed    & role & 1 & 2/100  & 0.02 & 0 \\
             & variable & none & 0 & 5/100  & 0.05 & 0 \\
             & variable & none & 1 & 1/100  & 0.01 & 0 \\
             & variable & role & 0 & 0/100  & 0.00 & 0 \\
             & variable & role & 1 & 0/100  & 0.00 & 0 \\
\addlinespace
Gemini       & fixed    & none & 0 & 8/100  & 0.08 & 0 \\
             & fixed    & none & 1 & 0/100  & 0.00 & 0 \\
             & \textbf{fixed} & \textbf{role} & \textbf{0} & \textbf{53/100} & \textbf{1.06} & \textbf{12} \\
             & fixed    & role & 1 & 6/100  & 0.06 & 0 \\
             & variable & none & 0 & 100/100 & 2.25 & 4 \\
             & variable & none & 1 & 0/100  & 0.00 & 0 \\
             & \textbf{variable} & \textbf{role} & \textbf{0} & \textbf{100/100} & \textbf{8.43} & \textbf{32} \\
             & variable & role & 1 & 50/100 & 1.16 & 2 \\
\addlinespace
GPT-4.1-mini & fixed    & none & 0 & 0/100  & 0.00 & 0 \\
             & fixed    & none & 1 & 0/100  & 0.00 & 0 \\
             & fixed    & role & 0 & 2/100  & 0.02 & 0 \\
             & fixed    & role & 1 & 0/100  & 0.00 & 0 \\
             & variable & none & 0 & 92/100 & 2.29 & 0 \\
             & variable & none & 1 & 1/100  & 0.01 & 0 \\
             & variable & role & 0 & 93/100 & 3.50 & 0 \\
             & variable & role & 1 & 0/100  & 0.00 & 0 \\
\addlinespace
GPT-4o-mini  & fixed    & none & 0 & 1/100  & 0.01 & 0 \\
             & fixed    & none & 1 & 0/100  & 0.00 & 0 \\
             & fixed    & role & 0 & 0/100  & 0.00 & 0 \\
             & fixed    & role & 1 & 1/100  & 0.01 & 0 \\
             & variable & none & 0 & 98/100 & 10.16 & 0 \\
             & variable & none & 1 & 0/100  & 0.00 & 0 \\
             & variable & role & 0 & 100/100 & 15.69 & 0 \\
             & variable & role & 1 & 1/100  & 0.01 & 0 \\
\bottomrule
\end{tabular}
}

\vspace{0.8em}
\noindent
Rebuttal Table 6 reduces those cells to simple effects on variable-arm participation, one row per
model.

\vspace{0.4em}
\noindent\textbf{Rebuttal Table 6.} Simple effects on participation in the variable arm, each factor evaluated with the other factor absent (pp). ROLE is the persona factor at RAT = 0; RAT is the rationality factor at ROLE = none; Mode is variable minus fixed at ROLE = none and RAT = 0.

\vspace{0.3em}
{\small
\begin{tabular}{lrrr}
\toprule
\textbf{Model} & \textbf{ROLE} & \textbf{RAT} & \textbf{Mode (variable $-$ fixed)} \\
\midrule
Claude-Haiku & $-5.0$ & $-4.0$ & $0.0$ \\
Gemini       & $0.0$  & $\mathbf{-100.0}$ & $\mathbf{+92.0}$ \\
GPT-4.1-mini & $+1.0$ & $\mathbf{-91.0}$  & $\mathbf{+92.0}$ \\
GPT-4o-mini  & $+2.0$ & $\mathbf{-98.0}$  & $\mathbf{+97.0}$ \\
\bottomrule
\end{tabular}
}

\vspace{0.8em}
\noindent\textbf{These are simple effects, not main effects.} Most cells sit at a floor or a
ceiling, and averaging over them would hide the persona $\times$ rationality interaction in
Gemini, where the persona moves instructed variable-arm participation from 0 to 50 of 100.

\vspace{0.3em}
\noindent\textbf{The headline.} Without the persona, the instruction reduces variable-arm
participation to 0--1 games in 100 in the three models that play most (Rebuttal
Table 6; Claude-Haiku was already floored). It is not a complete off switch: with the persona
present, Gemini still wagers in 50 of 100 variable games under the instruction and still goes
bankrupt twice.

\vspace{0.3em}
\noindent\textbf{However, this factorial cannot yet decide between your hypothesis and a
simpler one.} Our instruction combines information, evaluation and permission in one sentence,
so the design cannot separate ``the model never represented stopping as optimal'' from plain
normative instruction following; a decomposition arm would, and we have not run it. Nor did we
test the second half of your question: the instruction states expected value inside the casino
narrative, which no condition removes, so we cannot say whether the behaviour survives a
neutral decision-theoretic presentation.

\vspace{0.3em}
\noindent\textbf{Four limits bound this result.}
\textit{(i)} Bankruptcy, the pre-registered primary endpoint, is zero in most completed cells
but not all: Gemini's cells contain 50 bankruptcies, so the endpoint is informative there, and
fallback counts are zero in every completed cell.
\textit{(ii)} Claude-Haiku is at the stopping floor in every cell, so its null results do not
test an autonomy effect, and its small fixed-arm counts carry the parser annotation from
Weakness 3.
\textit{(iii)} The 12 open-weight cells are still running --- LLaMA, which carries the headline
behavioural effect, is among them --- and we report them by 3~August.
\textit{(iv)} GPT-4.1-mini carries a standing system message in every cell (``You are a
cautious, rational decision maker\ldots''), and the factorial prepends its factors to the user
prompt only, so its instruction-absent cells are not instruction-free. Its within-model
contrasts remain internally valid (the message is constant), its participation result is
conservative (92 of 100 variable games wagered despite the standing caution), and its
rationality effect should not be compared across vendors; a genuinely instruction-free baseline
cell is registered (Rebuttal Table 9).

\vspace{0.3em}
\noindent
One short instruction removes participation in the three models that play most. That is the
safety-relevant result your question produced, and the limits above bound it. We intend to
reflect this in the camera-ready.

\plan{합리성 지시 요인설계를 본문 결과로 편입하고, 지시가 정보·평가·허가를 묶고 있다는 점과 카지노 프레이밍 미제거를 한계로 명시하며, 분해 arm과 GPT-4.1-mini 지시-부재 기준셀을 2차 응답으로 등록한다.}

% ---------------------------------------------------------------------
\subsubsection*{\textbf{[Question 4]} Why SAE top-200 plus Ridge rather than logit controls or a behavioural-state baseline?}
% ---------------------------------------------------------------------

\textbf{Verdict: accepted in part. We have not run the behavioural-state baseline you name}, so
we cannot yet say whether the readout adds anything over quantities we already observe.

\vspace{0.3em}
\noindent\textbf{The original motivation was legibility, not statistical performance.} Selecting
the top-$k$ sparse-autoencoder (SAE) features and fitting ridge regression lets the recovered
direction be expressed through individually inspectable units. That choice does not imply better
estimation than a dense probe, and the paper now says so.

\vspace{0.3em}
\noindent\textbf{Two directions are at issue, and only one passes our causal test.} The test
adds a direction to activations at several strengths and measures the change in betting against
twenty random directions pushed the same way. The behavioural axis --- fitted from the model's
betting, not a decoder --- passes: slope 0.0457 at $z = +4.45$, and removing it lowers betting
in both open-weight models ($-0.037$ and $-0.052$, $n = 200$ seed-matched pairs each). The SAE
readout direction, which carries the paper's monitoring claim, fails the same test, exceeding
the null band only at one strength in one direction. We claim causal control for the
behavioural axis, not for the readout, and every internal-state claim depends on that
distinction.

\vspace{0.3em}
\noindent\textbf{The control we ran} addresses estimator artefacts: a plain ridge probe on raw
activations with no SAE also moves behaviour (slope 0.0284, $z \approx +3$) and is nearly
orthogonal to the other axes, so the split between behaviour-moving and behaviour-neutral
directions is not an artefact of ridge versus mean-difference or top-$k$ versus dense features;
its source is unidentified.

\vspace{0.3em}
\noindent\textbf{The control we did not run is the one you name.} We have not fitted a baseline
of balance, round index and recent outcome, nor included choice-probability or logit features,
so we cannot say how much the readout adds beyond observable game state. This weakness remains
open; it is registered with an SAE-reconstruction control and stronger null bands (Rebuttal
Table 9), reported by 3~August.

\vspace{0.3em}
\noindent\textbf{However, the causal battery already excludes the confounding version of your
concern, though not the redundancy version.} The strength ladders are paired designs --- the
same seeds and strengths on the same states --- so game state cannot generate a difference
between the compared arms; a direction defined by account balance moves betting no more than a
random one ($z = +0.64$), and removing it raises Gemma's betting, the wrong sign for a balance
confound. Reviewer KuK5's Question 2 contains the full argument. We intend to reflect this in
the camera-ready.

\plan{SAE 상위-k + Ridge의 동기를 가독성으로 명시하고, readout 방향의 인과 실패와 행동축의 통과를 분리 서술하며, 중첩 행동상태 베이스라인을 미해결 약점으로 남긴 채 2차 응답으로 등록한다.}

% ---------------------------------------------------------------------
\subsubsection*{\textbf{[Limitations]} Artificial negative-EV tasks, and monitoring signals rather than mechanism}
% ---------------------------------------------------------------------

\textbf{Verdict: accepted on task realism, and adopted with one scope line on evidence status.}
The tasks are artificial negative-EV games and do not represent risk behaviour in real
financial, planning or tool-using agents; we now state this in the contribution claim itself,
restated in two layers so the evidence status of each is visible without the appendix.

\begin{enumerate}[leftmargin=1.8em]
\item \textbf{Behavioural.} At a matched per-round cap, discretion over stake size raises
sustained participation in five of six models, and bankruptcy, the pre-registered primary
endpoint, moves decisively in LLaMA, weakly in Gemini and not at all in the other four. These
results concern artificial negative-EV tasks under a declared prompt regime, claim no
generalisation to real-world agents, and describe a condition-dependent policy, not a trait.
\item \textbf{Internal.} Decision-time hidden states carry a readout of these contrasts, which
supports monitoring and identifies no mechanism: the sole causal result is a limited
intervention on a task-defined betting-aggression axis, and the SAE readout direction fails
that causal test. Nothing here establishes endogenous circuitry or transfers to deployments.
\end{enumerate}

\noindent
Your closing sentence --- exploratory monitoring signals rather than evidence of an underlying
mechanism --- is now the paper's own claim for the layer it describes, the internal-state
results. We do not extend that label to the behavioural layer, which is a
pre-registered-endpoint experimental effect rather than an exploratory signal. We intend to
reflect this in the camera-ready.

\plan{기여 진술을 행동층·내부층 2단으로 재진술하고, 인공 부기대값 과제라는 범위 제한을 기여 문장 안에 넣으며, 내부 상태 결과는 탐색적 모니터링 신호로만 주장한다.}
